\documentclass{article}

\usepackage{amsmath}

\title{Ch. 24 Classwork Solutions}
\author{Principles of Macroeconomics}
\date{\today}

\begin{document}

\maketitle

\section*{Quick Quiz}

\begin{enumerate}

\setcounter{enumi}{10}
\item (c) -- GDP omits activity that takes place outside of markets like household production, and here Mr. Keating leaves the labor market to engage in household production.

\item (c) -- GDP doesn't account for the quality of the environment. 

\end{enumerate}

\section*{Questions for Review}

\begin{enumerate}

\setcounter{enumi}{7}
\item GDP is a measure of national income. That means higher GDP countries can afford higher levels of healthcare, education, etc. 

Implementing a seven day work week would increase GDP because we would produce more, but it wouldn't necessarily increase our welfare because we enjoy having time for leisure.

\end{enumerate}

\section*{Problems \& Applications}

\begin{enumerate}

\setcounter{enumi}{3}
\item Including the value of used goods that are resold would increase GDP even though no additional production occurred. It would make GDP a less accurate representation of a country's production and therefore a less accurate reflection of that country's ability to provide for itself. 

\item 

	\begin{enumerate}
	
	\item 
		\underline{2023} \\
		Nominal GDP $= 1(100) + 2(50) = 200$ \\
		Real GDP = Nominal GDP = 200 \\
		GDP Deflator = 100
	
		\underline{2024} \\
		Nominal GDP $= 1(200) + 2(100) = 400$ \\
		Real GDP $= 1(200) + 2(100) = 400$ \\
		GDP Deflator = 100
	
		\underline{2025} \\
		Nominal GDP $= 2(200) + 4(100) = 800$ \\
		Real GDP $= 1(200) + 2(100) = 400$ \\
		GDP Deflator $= \frac{800}{400} \cdot 100 = 200$
		
	\item
		\underline{2024} \\
		\begin{gather*}
		\% \Delta GDP_{Nom} = \frac{400 - 200}{200} \cdot 100 = 100 \\
		\% \Delta GDP_{Real} = \frac{400 - 200}{200} \cdot 100 = 100 \\
		\% \Delta GDP_{Def} = \frac{100 - 100}{100} \cdot 100 = 0
		\end{gather*}
		
		\underline{2025} \\
		\begin{gather*}
		\% \Delta GDP_{Nom} = \frac{800 - 400}{400} \cdot 100 = 100 \\
		\% \Delta GDP_{Real} = \frac{400 - 400}{400} \cdot 100 = 0 \\
		\% \Delta GDP_{Def} = \frac{200 - 100}{100} \cdot 100 = 100
		\end{gather*}
		
		In 2024, the quantities changed, but the prices didn't, so nominal GDP and real GDP increased while the GDP deflator remained constant. In 2025, the prices changed, but the quantities didn't, so nominal GDP and the GDP deflator increased while real GDP remained constant. 
		
	\item Economic well-being increased more in 2024 because in 2024 production increased. As a result, the country can produce more for itself.
		
	\end{enumerate}
	
\item 

	\begin{enumerate}
	
	\item Year 1: $3(4) = 12$ \\
		Year 2: $4(5) = 20$ \\
		Year 3: $5(6) = 30$
	
	\item I'll use year 1 as my base year. \\
		Year 1: $3(4) = 12$ \\
		Year 2: $4(4) = 16$ \\
		Year 3: $5(4) = 20$ \\
		
	\item Year 1: 100 \\
		Year 2: $20 / 16 * 100 = 125$ \\
		Year 3: $30 / 20 * 100 = 150$
		
	\item $\frac{20 - 16}{16} \cdot 100 = 25 \%$
	
	\item $\frac{150 - 125}{125} \cdot 100 = 20\%$
	
	\item I could have calculated the growth rate of the number of bars for (d), and I could have calculated the growth rate of the price of bars for (e). 
	
	\end{enumerate}
	
\item

	\begin{enumerate}
	
	\item $100 \left[ \left( \frac{21,141}{10,287} \right)^{\frac{1}{20}} -1 \right] = 3.67 \%$
	
	\item $100 \left[ \left( \frac{113.6}{78.1} \right)^{\frac{1}{20}} - 1 \right] = 1.89 \%$
	
	\item $\frac{10,287}{\frac{78.1}{100}} = 13,171.57$
	
	\item $\frac{21,141}{\frac{113.5}{100}} = 18,610.04$	
	
	\item $100 \left[ \left( \frac{18610.04}{13171.57} \right)^{\frac{1}{20}} - 1 \right] = 1.74$
	
	\item It must be higher because it also includes growth in prices.
	
	\end{enumerate}
	
\setcounter{enumi}{9}
\item If much of the production in India takes place outside of markets and if many of the trades are unrecorded, then India's reported real GDP per person will be substantially lower than the country's actual production per person. As a result, the gap in economic well-being between the United States and India will be overstated.

\item

	\begin{enumerate}
	
	\item It increased GDP because many women's labor transitioned from unmeasured household production to measured market production.
	
	\item The increase in GDP would be smaller because the increase in market production would be partially offset by the decrease in household production.
	
	\item The quality of childcare might have changed. If mothers joined the labor force and fathers didn't exit, parents must have outsourced child care to other family members, nannies, or day care. Presumably, the quality of childcare varies across those three options and a stay-at-home parent.
	
	\end{enumerate}

\end{enumerate}

\end{document}